\section{Related Work}
\label{sec:related}

PEM sits at the intersection of five threads: semi-supervised medical image segmentation, classical entropy-minimization SSL, test-time adaptation, post-hoc calibration, and post-hoc checkpoint refinement. We discuss each in turn and make the contrast with PEM explicit, because the differences are regime-level and not always apparent from the loss functions alone.

\paragraph{Semi-supervised medical image segmentation.}
Two mechanisms dominate the medical SSL literature: consistency regularization --- where a student network is trained to agree with a teacher or with itself under perturbations of the input, the model, or the task --- and pseudo-labeling, where high-confidence student predictions are converted to hard targets and folded back into the supervised loss. The mean-teacher framework~\cite{tarvainen2017meanteacher} instantiates the first mechanism with an exponential-moving-average teacher; UA-MT~\cite{yu2019uamt} adds MC-dropout-based uncertainty filtering so that the consistency loss only acts on voxels the teacher is confident about; DTC~\cite{luo2022dtc} enforces consistency between the segmentation head and a level-set head to force geometrically coherent predictions; MC-Net+~\cite{wu2022mcnet} uses cross-pseudo-supervision between multiple decoders; CPS~\cite{chen2021cps} applies the same idea to two networks initialized differently. BCP~\cite{bai2023bcp} introduces bidirectional copy-paste between labeled and unlabeled volumes, mixing them at the pixel level during training so that every forward pass exposes the network to both distributions. MOST~\cite{liu2024most} adds multi-formation soft masking to the consistency objective, and DyCON~\cite{assefa2025dycon} couples dynamic uncertainty weights with a contrastive loss. What all of these methods share, despite their mechanistic differences, is the training-time contract: a composite loss that couples labeled and unlabeled data is optimized to convergence, and the final checkpoint is then released as the end product. The unlabeled data is not revisited. PEM is the first method in this family to ask what the discarded unlabeled data still contains at convergence, and to turn the answer into a refinement step.

\paragraph{Classical entropy minimization and its modern descendants.}
Grandvalet and Bengio~\cite{grandvalet2004entropy} introduced entropy minimization as a semi-supervised regularizer, adding $-\lambda H(Y\mid X)$ to the supervised loss on unlabeled data and showing that the resulting classifier implements the low-density separation principle associated with the cluster assumption~\cite{chapelle2006ssl}. The classifier that minimizes conditional entropy pushes its decision boundary into low-density regions of input space, so the same loss acts as a margin-maximization regularizer when classes form clusters. Pseudo-Label~\cite{lee2013pseudolabel} is the hard-label limit of the same idea: entropy minimization with one-hot targets. Modern SSL methods such as MixMatch~\cite{berthelot2019mixmatch}, ReMixMatch~\cite{berthelot2020remixmatch}, and FixMatch~\cite{sohn2020fixmatch} incorporate entropy sharpening of weak-augmentation predictions as a way to generate pseudo-targets for strong-augmentation consistency, and FlexMatch~\cite{zhang2021flexmatch} couples this with a curriculum of class-specific confidence thresholds. In every case, the entropy term is active during training, coupled to a supervised or consistency loss that prevents it from collapsing the posterior to a single class, and requires careful balancing of its strength over time. The entropy loss we use is the same formal object as in Grandvalet and Bengio; what is new is the regime. PEM applies the loss to a checkpoint that has already converged under a different training objective, and the residual entropy it sees is structurally concentrated at decision boundaries rather than diffuse across the input. The gradient signal is automatically well-localized, and none of the stabilizing machinery required by training-time entropy minimization --- class-balance regularizers, warmup schedules, threshold curricula, auxiliary losses --- is needed. To the best of our knowledge, this is a regime that has not been isolated in the prior entropy-minimization literature, and whose properties have not been characterized.

\paragraph{Test-time adaptation.}
The methodologically closest body of work is test-time adaptation (TTA), in which a pre-trained model is adjusted on test-time inputs to handle distribution shift. TENT~\cite{wang2021tent} minimizes prediction entropy on the test stream by updating only the BatchNorm affine parameters; the loss is identical to the one PEM uses. TTT~\cite{sun2020ttt} and TTT++~\cite{liu2021tttplusplus} use auxiliary self-supervised objectives at test time; CoTTA~\cite{wang2022cotta} and EATA~\cite{niu2022eata} address continual adaptation under non-stationary streams. The similarity to PEM is real and worth making explicit: both TENT and PEM minimize the same loss on unlabeled data after training. The differences are equally real. First, TTA assumes a distribution shift between the training set and the test stream, and its purpose is to recover performance lost to that shift. PEM makes no distribution-shift assumption: it is applied to unlabeled data drawn from the same distribution as the labeled training set, and its purpose is to recover the signal that the base SSL training left on the table at convergence. Second, TTA typically operates on per-sample or per-batch streams and does not produce a deployable checkpoint in the usual sense --- the model's parameters depend on the test stream seen so far. PEM produces a single deployable checkpoint in a one-time fine-tune, evaluated with a standard frozen-weights protocol. Third, and most importantly, TTA's evaluation protocol is constructed around distribution shift: reported numbers compare the test-time-adapted model against the shift-degraded baseline, not against the non-shifted upper bound. PEM's evaluation compares the fine-tuned checkpoint against the best available SSL training method on the same distribution, using the same test set as the original benchmark. The two methods share a loss but solve different problems and are measured against different baselines; results from one do not directly transfer to the other. We view TENT as a natural starting point for thinking about PEM and as a method whose insights should be imported into the SSL-convergence regime, rather than as a competitor we are displacing.

\paragraph{Post-hoc calibration.}
Temperature scaling~\cite{guo2017calibration} rescales the output logits by a scalar $T$ chosen on a held-out set to minimize negative log-likelihood or expected calibration error. It is the canonical post-hoc fix for miscalibrated classifiers. Its key property is that it cannot move argmax: for any monotone rescaling of the logits, the voxel-wise argmax is preserved, and therefore so is the segmentation mask. As a consequence, temperature scaling cannot move Dice, Jaccard, HD95, or any other region-overlap metric. We include it as a controlled post-hoc baseline in our experiments precisely because it isolates \emph{recalibration} as a possible explanation for PEM's gains: if PEM were merely sharpening the posterior without changing the underlying decision surface, temperature scaling would produce the same improvement. Temperature scaling produces $|\Delta\mathrm{Dice}|\!\le\!0.03$ on all three of our configurations (Table~\ref{tab:main}), which rules this explanation out and confirms that PEM's improvement arises from parameter updates at decision-boundary voxels rather than from confidence rescaling.

\paragraph{Post-hoc model refinement.}
Several techniques refine a trained model after the fact without relying on additional labels. Stochastic weight averaging (SWA)~\cite{izmailov2018swa} averages weights across points in the training trajectory and produces flatter-minimum checkpoints at no additional training cost. Polyak averaging~\cite{polyak1992acceleration} and exponential-moving-average teachers produce closely related checkpoints during SSL training itself. Self-distillation~\cite{hinton2015distillation,allen2023self} distills a trained model into a fresh student using the teacher's soft predictions as targets. None of these methods uses entropy minimization, and all of them either require access to the training trajectory (SWA, EMA) or perform full supervised-style distillation with cross-entropy against the teacher (self-distillation), which is closer to pseudo-label fine-tuning than to PEM. The pseudo-label fine-tune baseline we include in Section~\ref{sec:experiments} is the direct post-hoc instantiation of self-distillation using hard argmax targets, and we show that it is outperformed by PEM on the two less-saturated configurations while becoming competitive in the saturation limit, in a manner that is quantitatively consistent with the theoretical argument of Section~\ref{sec:theory}.

\paragraph{Uncertainty estimation in segmentation.}
A distinct strand of work addresses \emph{measurement} of uncertainty in segmentation outputs --- Monte Carlo dropout~\cite{gal2016dropout}, deep ensembles~\cite{lakshminarayanan2017simple}, evidential networks~\cite{sensoy2018evidential}, and Bayesian SegNet~\cite{kendall2015bayesian}. These methods produce per-voxel uncertainty maps but do not, by themselves, reduce uncertainty or improve segmentation quality. Our diagnostic $\rho_f(\tau)$ uses the simplest possible per-voxel uncertainty proxy (softmax entropy) and aggregates it into a scalar characterization of where the uncertainty lives, rather than what its magnitude is. More sophisticated uncertainty estimators could be substituted for softmax entropy in the definition of $\rho_f$, and we expect the qualitative story to be preserved; we use softmax entropy for simplicity and because the PEM loss is computed against the same quantity, which makes the connection between the diagnostic and the loss tight.

\paragraph{Summary of positioning.}
PEM is not a new loss, a new architecture, or a new SSL training recipe. It is the observation that converged SSL checkpoints leave a specific kind of usable signal on the unlabeled data, a theoretical account of why that signal is in a uniquely favorable regime for entropy minimization, a scalar diagnostic that detects when this regime holds, a label-free fine-tuning step that exploits it, and an empirical validation that the resulting improvements are not explicable by recalibration, pseudo-labeling, consistency, or seed variance. The contribution is the regime shift and its consequences, not the loss.
